\documentclass[14pt,a4paper]{extarticle}
\usepackage[left=3cm,right=2cm,top=2cm,bottom=2cm,headheight=18pt]{geometry}
\usepackage{fontspec}
\setmainfont{Liberation Serif}
\setsansfont{Liberation Serif}
\setmonofont{Liberation Mono}
\usepackage{unicode-math}
\setmathfont{texgyretermes-math.otf}
\usepackage{amsmath,graphicx,booktabs,longtable,array,enumitem}
\usepackage[normalem]{ulem}
\usepackage{titlesec,fancyhdr}
\usepackage[hidelinks,unicode]{hyperref}
\usepackage{url}
\urlstyle{same}
\setlength{\parindent}{1.27cm}
\setlength{\parskip}{0pt}
\linespread{1}
\setlength{\emergencystretch}{2em}
\clubpenalty=10000
\widowpenalty=10000
\titleformat{\section}{\normalfont\bfseries}{\thesection.}{0.5em}{}
\titleformat{\subsection}{\normalfont\bfseries}{\thesubsection.}{0.5em}{}
\titlespacing*{\section}{0pt}{1.2em}{0.5em}
\titlespacing*{\subsection}{0pt}{0.9em}{0.3em}
\setlist{nosep,leftmargin=1.27cm}
\pagestyle{fancy}
\fancyhf{}
\renewcommand{\headrulewidth}{0pt}
\renewcommand{\footrulewidth}{0pt}
\fancyhead[C]{\ifnum\value{page}>1\thepage\fi}
% Liberation Serif is a disclosed Times-compatible substitute, not Times New Roman.
% Current statutory geometry/size baseline: O‘RQ-682 appendix §§43–51.
\hyphenpenalty=10000
\exhyphenpenalty=10000
\pretolerance=10000
\tolerance=10000
\setlength{\emergencystretch}{4em}
\renewcommand{\normalsize}{\fontsize{14bp}{16.8bp}\selectfont}
\normalsize
\begin{document}
\begin{center}\textbf{ПОЯСНИТЕЛЬНАЯ ЗАПИСКА}\end{center}
Научно-техническое обоснование авторского предложения о проверке расчётов Социального реестра. Редакция от 1 октября 2026 года; опорная дата расчётов — 30 сентября 2026 года. Основная русская редакция для подачи. Материал 2; материал 1 содержит нормативные положения, материал 3 — исследование с электронным вычислительным дополнением.

\section*{Основание, разработчики и цель}
Пункты 13–14 PF-193 предусматривают многомерную оценку с 1 января 2027 года и внесение проекта в Администрацию Президента до 1 декабря 2026 года. Официальные разработчики — Национальное агентство социальной защиты при Президенте Республики Узбекистан и Министерство экономики и финансов Республики Узбекистан. Записка относится к предложению Abduxoliq Akmal ugli Ashuraliyev и не подаётся от имени официального разработчика.

Цель — повторить утверждённый расчёт на одинаковых входах, проверить границы и исключения, отличить проблему данных от программного дефекта и объяснить основания индивидуального решения. Основное предложение содержит семь основных пунктов и приложение из двадцати восьми пунктов. Само оно не устанавливает баллы 2027 года или выплаты. Ожидается проверяемый и восстанавливаемый расчёт; количественный прогноз благосостояния, экономии или охвата не дан. Также дана краткая альтернативная редакция из пяти пунктов для замены полной формы, без совокупного применения.

Вид акта определяется основным проектом. Вариант, затрагивающий PF-258, требует изменения президентского акта. Правительственный акт и программная настройка не заменяют вышестоящий критерий.

\section*{Предмет регулирования и структура предлагаемого акта}
Основное решение — обязательная проверка исполнения уже утверждённых правил, а не выбор новой модели нуждаемости. Семь основных пунктов устанавливают утверждение порядка, единое описание методики, воспроизводимость, организацию подготовки по статье 22, согласование связанных актов, ресурсный расчёт и готовность до применения. Приложение содержит семь глав: общие положения (пункты 1–3); описание и версии (4–7); качество данных и сбои (8–11); проверка расчётов (12–17); основания и пересмотр решений (18–20); исследование, отчёт и контроль (21–24); приёмка, исправление и хранение (25–28). Нормативные требования сосредоточены в материале 1; эта записка объясняет их назначение и исполнение, материал 3 содержит подробные доказательства.

Проект относится к подзаконному регулированию исполнения поручения PF-193. Требование пункта 84 Единой методики рассмотреть перенос подзаконных отношений на уровень закона относится к подготовке законопроекта; самостоятельный проект закона здесь не предлагается. Компетентный орган выбирает форму основного акта. Новое правило о праве на помощь или размере выплаты устанавливается только актом надлежащей юридической силы. Проверка расчёта не заменяет этого полномочия.

\section*{Необходимость каждого блока требований}
Действующие обязанности правильной работы и своевременного предоставления сведений по пункту 46 приложения 1 к постановлению № 35 являются отправной точкой. Предлагаемый порядок делает их исполнение проверяемым: для конкретной версии можно установить, какое правовое правило применено, какими исходными значениями оно обеспечено, совпали ли контрольные и фактические компоненты и кто разрешил обнаруженное расхождение. Обязанности заявителя, социального работника и «махаллинской семёрки» по тому же пункту сохраняются; проверка не переносит ответственность государственного источника на семью.

Паспорта критериев и датированные версии нужны потому, что одинаковое число при разных единицах, периодах или правилах включения не означает одинаковый вход. Раздельные состояния качества нужны, чтобы отсутствие записи не превращалось в подтверждённый нулевой доход, а задержка получения — в более достоверное значение. Контроль исключений и порядка их применения нужен, потому что правильная общая формула может дать неверное решение при пропущенном специальном условии. Компонентное сравнение нужно, поскольку две ошибки могут компенсироваться в итоговой сумме. Индивидуальное объяснение и связанное исправление нужны для установления конкретного основания решения и законного пересмотра, а не только для статистического отчёта.

Порядок содержит минимальные обязательные сведения и действия. Внутренняя форма паспорта, протокола и журнала может определяться Агентством без изменения критерия, срока обращения или основания отказа. Если разработчик предъявляет действующее равноценное требование и фактически обеспеченный контроль, соответствующий пункт интегрируется в него. Простая ссылка на наличие системы или общего контроля не разрешает отдельно доказанные пробелы единиц, границ, исключений и объяснения.

\section*{Выбор правового инструмента по каждому пункту}
Статья 20 Закона ЗРУ-682 требует проверить достаточность действующих механизмов и административных полномочий. Поручение УП-193 устанавливает задачу подготовки, но не доказывает необходимость новой нормы для каждого контроля. Здесь Н означает предлагаемую публичную обязанность или полномочие, требующее компетентного нормативного принятия; С — согласование существующей обязанности, которую при доказанной эквивалентности объединяют, а не дублируют; В — внутреннее исполнение в существующих полномочиях. Смешанные записи отделяют публичный минимум от исполнения. Это обоснованная классификация автора, а не официальное установление отсутствия полномочий Агентства. Действующие механизмы включают пункты 45–50 приложения 1 к постановлению № 35 и статьи 11 и 24 Закона о персональных данных.

\begingroup\fontsize{11bp}{13.2bp}\selectfont
\begin{longtable}{p{0.12\linewidth} p{0.13\linewidth} p{\dimexpr0.75\linewidth-6\tabcolsep\relax}}
\toprule Пункт & Выбор & Основание и решение\\\midrule\endhead
Осн. 1 & Н/В & Принять публичный минимум; формы делегировать. Переименование существующего контроля не требует отдельного акта.\\
Осн. 2 & Н/С & Критерии требуют компетентного утверждения; описание не меняет их самостоятельно.\\
Осн. 3 & С/В & Правильная работа системы и законные сроки уже обязательны; проверку организовать существующим полномочием.\\
Осн. 4 & С/В & Организация по статье 22 и внутренний регламент; нет новой внешней квалификации или обязанности гражданина.\\
Осн. 5 & С/В & Согласовать правила и процедуры; нет резервного внесения или изменения начала.\\
Осн. 6 & С/В & Дополнительные задачи, покрытие и реальные средства; нет повторной закупки.\\
Осн. 7 & С/В & Записи выпуска и согласования; новый отчёт не условие запуска.\\
Прил. 1 & В & Область и связанный набор контроля организуют существующую ответственность.\\
2 & В & Термины обеспечивают единое исполнение, не самостоятельное право на помощь.\\
3 & С & Категория отделена от условий помощи; отраслевые нормы сохраняются.\\
4 & В & Детали паспорта и источников переносятся во внутренний регламент.\\
5 & Н/В & Публичные критерии разрешают равенство и приоритет; программа проверяет исполнение.\\
6 & В/С & Версии и правовые даты обеспечивают восстановление по существующим обязанностям.\\
7 & С/В & Утверждающий акт и законное описание на существующем ресурсе; внутренняя версия.\\
8 & В/С & История источников обеспечивает основания и точность; схема записи внутренняя.\\
9 & Н/С & Отсутствие данных не новое неблагоприятное основание; законное вменение сохраняется.\\
10 & Н/С/В & Компетентное уточнение, сроки и отсутствие нового отказа; формы направления делегируются.\\
11 & С/В & Учёт существующего полномочия при сбое без расширения его области.\\
12 & В & Независимые контроли, охват и сверка состава исполняют требование правильной работы.\\
13 & В & Предварительный план управляет трактовкой результатов, не обязанностью заявителя.\\
14 & В & Числитель и знаменатель относятся к отчёту, не нуждаемости.\\
15 & В & Эмпирическая точность требует доказательств; исследование не условие запуска.\\
16 & В & Доказательства протокола и формы относятся к внутреннему регламенту.\\
17 & С/В & Исправление и выявление решений; индивидуальные изменения требуют законных полномочий.\\
18 & С/В & Доступные основания сохраняются; формат технической истории внутренний.\\
19 & Н/С & Действующее объяснение согласуется; доступный расчёт дополняет SMS.\\
20 & С/В & Законные сроки возражения и исправления; нет обязательного досудебного этапа.\\
21 & С/В & Действующие основания обработки и специальные условия; разрешение исследования внутреннее.\\
22 & С/В & Защищённое исполнение и раскрытие; нет внешнего доступа к персональным данным.\\
23 & С/В & Внутренний протокол и законный доступ/раскрытие; нет новой периодической отчётности.\\
24 & С & Существующий надзор; нет нового органа или основания уменьшения выплаты.\\
25 & В/С & Действующая запись выпуска; внешние обязательные заключения сохраняются.\\
26 & С/В & Исправить/проверить программу по действующему правилу; нет ручного резерва или истёкшего критерия.\\
27 & С & Даты и законный пересмотр; нет освобождения от переходного долга.\\
28 & С/В & Законное хранение и уничтожение; расписание категорий внутреннее.\\
\bottomrule\end{longtable}\endgroup

\section*{Необходимость нормативного содержания и границы технического решения}
\textbf{Предмет необходимости.} Пункты 13–14 \href{https://lex.uz/docs/-8472526}{PF-193} устанавливают новое многомерное распределение по баллам и поручают нормативный проект. Это прямое основание работы над нормативным содержанием 2027 года. Оно не доказывает обязательность отдельного акта только о тестировании или всех 28 авторских пунктов. Предлагается включить необходимое содержание в уже порученный акт: пяти пунктов краткой альтернативы материала 1 достаточно как формы предложения. Правила баллов, коэффициенты и все параметры 2027 года здесь не выдуманы; их ещё должен законно определить разработчик.

\textbf{Почему техническое решение не заменяет правового.} По статьям 3, 14, 17, 18 и 37–40 \href{https://lex.uz/docs/-5378966}{ЗРУ-682} значение имеют нормативный характер, компетенция, юридическая сила, установленное оформление, регистрация в применимых случаях и официальная публикация. Если выбор порога, равенства, периода, допустимого входа или вменения сам устанавливает условие признания либо отказа, его нельзя оправдать только успешным тестом. Тест отвечает, исполняется ли выбранное правило; компетентный акт отвечает, почему именно это правило законно применяется к семье. Это вывод автора из указанных норм, а не утверждение, что каждый программный параметр подлежит отдельному нормативному утверждению. Действующая законная делегация и техническое исполнение уже определённого условия сохраняются.

\textbf{Разграничение документов.} Пункты 94–95 \href{https://lex.uz/docs/-7766987}{Правил, зарегистрированных за № 3686 от 9 октября 2025 года}, охватывают ведомственные нормативные акты, затрагивающие права и гарантии, имеющие межведомственный характер либо обязательные для организаций вне системы принявшего органа. Пункт 96(b) исключает из регистрации документы исполнения ранее установленного порядка без новых правовых норм; 96(v)–(j) предусматривает технические исключения, включая отнесение иных документов Министерством юстиции. Поэтому название «техническое требование» само по себе не решает правовую природу содержания. Обратное тоже верно: подлинно технический документ не превращается в нормативный только из-за этого предложения. Эти правила относятся к ведомственной форме и не объявляют регистрацию условием президентского или правительственного акта. Проверена редакция с изменениями от 17 сентября 2026 года; пункт 20 прямо требует изучить необходимость и альтернативы.

\textbf{Два независимых условия эквивалентности.} Для правового покрытия указываются конкретный пункт, компетентный орган, юридическая сила, область семей/случаев, даты применения и обязательное последствие, а для ведомственной формы — применимая регистрация и публикация. Для технического покрытия указываются входы, даты, исключения, независимо обоснованный ответ, версия и выполненный результат. Охват всех авторских тестов устраняет необходимость повторить тесты. Он устраняет необходимость дополнительной нормы лишь когда правовое содержание также полностью обеспечено. Общий договор, аудит, приказ о составе группы или протокол приёмки не принимается за доказательство отсутствующего критерия. Если оба покрытия подтверждены, дополнительная норма по этому вопросу не предлагается.

\textbf{Как разрешается ограничение области или толкование.} Компетентное толкование существующего условия либо документированная область исходных данных могут снять конкретный контрпример. Указываются источник, полномочие, период, единицы и фактическая семантика полей; математический вывод тогда ограничивается этой областью. Текущее программное отбрасывание значения само по себе не доказывает, что семья юридически не вправе представить соответствующий факт. Новое ограничение признания не вводится под видом доказательства невозможного входа. По статье 55 ЗРУ-682 толкование не может менять, дополнять или исправлять акт. Выясненный смысл существующей нормы требует согласованного исполнения, а новый выбор правового условия — компетентного изменения. Закрытый вопрос не сохраняется как искусственный аргумент новой нормы.

\textbf{Прямо существующее временное полномочие.} Пункт 48 приложения 1 к \href{https://lex.uz/docs/-8027513}{постановлению № 35} охватывает технические неисправности, перерывы и неполное получение электронных сведений и разрешает Агентству временно прекратить обмен в отдельных системах до полного восстановления и изменить алгоритм определения категории. Поэтому не заявляется, что любого временного изменения алгоритма требует новый акт. Краткий пункт 4 и полный пункт 11 сохраняют это полномочие и предлагают прослеживаемость его конкретного применения. Объём его распространения на новую методику согласовывается разработчиком; наличие пункта 48 не объявляется достаточным правовым определением всей постоянной системы баллов 2027 года.

\textbf{Стоимость и порядок выбора.} Сначала в порученном проекте определяются законные условия оценки и сохранённые гарантии; одновременно сопоставляются существующие тесты и договоры. Затем оценивается только недостающая интеграция. Краткая редакция не учреждает новый поток обработки заявлений или национальную ручную службу. Юридическая работа, исполнение тестов и изменение интерфейса всё же требуют времени; ни отсутствие отдельного ассигнования, ни обещание внутреннего исполнения не доказывает нулевой стоимости. Неизвестные часы, договорное покрытие, мощность и финансирование здесь не подтверждены. Обоснованная ресурсная трудность может изменить объём технической формы и распределение работы, но не является правовым основанием позволить коду установить отсутствующее условие помощи.

\textbf{Критерий принятого решения.} Допускаются другой порядок пунктов, короче сформулированная норма, отсылка к уже действующему акту и надлежаще делегированный ведомственный нормативный акт. Для каждой правовой задачи сопоставляются сохранённое обязательное последствие, источник, даты и компетенция; для технической — исполнитель и эквивалентное доказательство. Такое сопоставление запрашивается для обсуждения, а не объявляется новой обязательной формой государственного ответа. Включение в опубликованный компетентный акт равнозначной нормы является нормативным результатом и при отсутствии отдельного приложения. Включение случаев только в технические требования является техническим результатом; оно не закрывает доказанно неурегулированное нормативное содержание.

\begingroup\fontsize{11bp}{13.2bp}\selectfont
\begin{longtable}{p{0.26\linewidth} p{0.39\linewidth} p{\dimexpr0.35\linewidth-6\tabcolsep\relax}}
\toprule Возражение & Достаточное подтверждение & Решение по предложению\\\midrule\endhead
Действующий контроль уже покрывает случаи & Применимая норма и независимо обоснованные выполненные проверки той же области и дат & Повторный тест исключается; норма объединяется только при подтверждённом правовом покрытии.\\
Контрпример вне области либо уже разрешён & Авторитетный источник ограничения или компетентное толкование, применимая дата и семантика входа & Случай ограничивается или снимается; новое правовое ограничение требует компетентного акта.\\
Дополнительная работа & Конкретная непокрытая задача, часы, договор, средства и при изменении обработки её нагрузка & Объём технической интеграции соразмеряется; правовое условие не заменяется настройкой.\\
Нужна другая редакция & Та же компетенция, обязательное последствие, область, даты и исполнимый канал & Краткая или эквивалентная нормативная редакция принимается как результат; точный авторский текст не является целью сам по себе.\\
\bottomrule\end{longtable}\endgroup

\textbf{Минимальный вариант принятия.} В порученный акт включаются необходимые правовые условия и специфичные для оценки гарантии; материал 1 предлагает краткую альтернативу из пяти пунктов. Полные семь пунктов и приложение — другая форма, не совокупное требование. Приёмочные тесты и подробные формы используют существующие полномочия и равноценные записи. Техническое включение научных случаев достаточно для инженерной задачи, когда правовое покрытие отдельно подтверждено. Оно не заменяет принятие отсутствующего правового условия. Сокращение либо отсылка допустимы при сохранении компетенции, обязательного последствия, области, дат, доступа и законных сроков.

\section*{Доказательства и поддерживаемая ими узкая мера}
\begingroup\fontsize{11bp}{13.2bp}\selectfont
\begin{longtable}{p{0.30\linewidth} p{0.35\linewidth} p{\dimexpr0.35\linewidth-6\tabcolsep\relax}}
\toprule Выполненное доказательство & Обоснованная узкая мера & Что не установлено\\\midrule\endhead
Контроли равенства дохода и трудной ситуации & Согласовать равенство, приоритет исключения и законные ожидаемые ответы. & Частота ошибок системы или преимущество включающей политики границы.\\
620 отрицательных результатов площади в 4\,105 построенных строках & Подтвердить геометрию и область, решить правовую формулу до программирования; проверить выбранное правило. & 620 пострадавших семей, национальная частота или выгода нулевого ограничения.\\
Официальный пример: 432\,000 сумов & Проверить единицы, делитель семьи и период по правовому примеру. & Общая правильность действующей системы.\\
Контроли налога, скота и условных переводов & Сохранить ветви, единицы и датированные параметры в проверках. & Реализация Агентства, полная спецификация 2027 года или качество реальных данных.\\
37 регрессионных тестов и проверка фиксированного состава & Воспроизводимый инструмент без пропуска или дублирования построенных случаев. & Приёмка системы, доступные работники Агентства или экономия.\\
\bottomrule\end{longtable}\endgroup
Контроли устанавливают вопросы для разрешения, не измеренную дополнительную пользу относительно реализации Агентства. Сопоставляются спецификация, независимо обоснованные ответы, записи выпуска, объяснение и обязательные технические/защитные заключения. Равноценный контроль используется повторно; дополнительна только выявленная непокрытая задача. Данные семей не выходят из Агентства. Авторитетная область может исключить сконструированный вход. Оценка благосостояния имеет отдельную область. Новая публичная статистическая отчётность не предлагается.

\section*{Внутренняя техническая спецификация вне нормативного приложения}
Это авторское задание для внутреннего регламента основного пункта 4, а не дополнительный публичный критерий, обязанность заявителя или полномочие менять законодательство.

\textbf{Паспорта критерия и источника.} Сохраняются постоянный идентификатор, правовое правило, тип и область, собственник и идентификатор записи, единица, обязательность, границы периода, обновление, преобразование, дублирование и пересечение, выбор при конфликте, законное вменение, порядок операций и промежуточная точность. Перевод и порог имеют один период и единицу; состав семьи и меняющаяся оплата труда — применимую дату. Поступление записи не является датой факта. Сохраняются конфликтующие значения и законная причина выбора. Исправление работником имеет автора, время и полномочие.

\textbf{Реестр контроля.} Каждое применимое условие, граница, исключение, приоритет и достижимое взаимодействие связываются со случаями. Для случая указываются идентификатор, входы, предпосылки, дата действия, законно выведенные промежуточные и окончательные ответы, вывод и независимый проверяющий. Сверяются запланированные и выполненные идентификаторы и кратность; исключение и неприменимость объясняются. Пустой состав и неопределённый знаменатель показываются явно. Изменение правила или источника требует проверки затронутых и взаимодействующих путей; при неопределённой области влияния проверяется вся применимая цепочка. Общий вывод программ не является независимым правовым ответом.

\textbf{Протокол и исправление.} Отдельно сохраняются идентичность методики, программы и конфигурации, период и состав данных, область и реестр, ожидаемые и фактические компоненты, исполненные и исключённые количества, различия, влияние, ответственный за исправление и повторный тест. Отделяются соответствие правилу, качество данных и полнота оснований. Дефект связывается с версиями, датами и заявлениями, а не только обнаружившей его жалобой. Уточнение связывает первоначальное поступление, источник либо получателя, законный срок, ответ, новый расчёт и уведомление. Защищённые записи остаются в разрешённом хранилище, публичная выписка содержит только проверенные на раскрытие агрегаты.

\section*{Соответствие при запуске без нового перехода прав на помощь}
Пункт 26 теперь касается исправления, повторной проверки и законного управления программным выпуском. Второй службы ручного решения он не требует. Прежний выпуск используется только при исполнении законно применимого к решению правила и проверке относящихся к нему контролей и пригодности к работе. Выпуск с заменённым правилом дохода не является законным резервом лишь потому, что прошёл прежние тесты. Меры сбоя ограничены пунктом 48 приложения 1 к постановлению № 35; иной дефект не расширяет это полномочие.

Поданный нормативный текст не содержит резервного перехода Президента с сохранением методики 2026 года, включающей границей дохода или освобождением от переходного долга. Эти решения требуют отдельных доказательств правовых и бюджетных последствий и не нужны для узкого вклада в проверку соответствия. Варианты равенства и нижней границы земельной формулы ниже остаются условными исследовательскими вариантами, не запрашиваемыми условиями принятия. Исключение резерва не означает, что программа устраняет отсутствие нормы: разработчики завершают законную таблицу решений и разрешают отсутствие или противоречие правила до первого применения. По статье 55 O‘RQ-682 компетентное официальное толкование объясняет подзаконную норму, но не вносит изменение. Критерий либо начало действия меняет только компетентный принимающий орган. Назначенная помощь, ожидающие заявления и взыскание регулируются применимым законодательством; промежуточное право этим предложением не создаётся.

\section*{Действующие контроли, операционные сведения и дополнительная задача}
Следующие официальные материалы открыты и проверены 1 октября 2026 года. Правовая обязанность отличается от доказательства работы услуги либо мероприятия. Полного реестра проверок программного выпуска Агентства и измеренного ресурсного расчёта эти материалы не содержат.
\begingroup\fontsize{11bp}{13.2bp}\selectfont
\begin{longtable}{p{0.30\linewidth} p{0.33\linewidth} p{\dimexpr0.37\linewidth-6\tabcolsep\relax}}
\toprule Официальное доказательство & Установленная область & Узкое включение и предел\\\midrule\endhead
Постановление № 35, приложение 1, пункты 45–50 & Защита данных, ответственные за систему и источники, ограниченное полномочие при сбое, контроль и споры. & Используются эти исполнители и записи. Наличие независимого ожидаемого ответа для каждого пути категории не доказано.\\
\href{https://lex.uz/docs/-6200851}{Постановление № 516}, основной пункт 6$^{1}$, приложение 1, пункты 13 и 44 & Обязательные положительные технические и защитные заключения до внедрения; совместные тесты подключения; доступ к количеству запросов, периодам и статистике. & Проверка правила и ожидаемого ответа включается в применимые технические требования и приёмку. Успех подключения не доказывает арифметику категории; требуемые внешние заключения не заменяются.\\
\href{https://lex.uz/docs/-8076376}{Постановление № 93 от 5 марта 2026 года}, основной пункт 1 и пункт 3 приложения & Функция IT-аудита Министерства цифровых технологий, включая процессы инфраструктуры и управления изменениями. & Используются относящиеся результаты аудита, не создаётся повторная инспекция. Результат конкретного аудита скоринга Агентства актом не опубликован.\\
\href{https://oldmy.gov.uz/oz/service/587}{Официальный портал, услуга 587} & Действующий канал заявления и результата, отсутствие требуемых документов, бесплатность, указанный срок в один месяц, опубликованные количества статусов. & Используется существующий канал. Просмотренная страница показывает всего 640\,876, в процессе 16\,564; период накопления, возраст ожидающих решений и время обработки не указаны. Это не годовое поступление, число просрочек или доступная мощность.\\
\href{https://gov.uz/oz/news/view/194551}{Сообщение Правительства от 17 июля 2026 года} & Сообщено о 2\,000 необоснованных отказах в 184 махаллях за месяц и росте обращений Президенту по реестру с 30\,000 до 46\,000; поручен прокурорский контроль. & Обосновывает доступ к основаниям дела и пересмотр по конкретной причине. Программная причина, нынешняя частота и эффект предложения не установлены.\\
\href{https://gov.uz/oz/urganchshahar/news/view/134892}{Сообщение хокимията Ургенча от 19 февраля 2026 года} & Проведён практический семинар о законном включении/исключении и достоверности данных. & При пригодности используются существующие каналы обучения; местное мероприятие не доказывает общереспубликанскую готовность специалистов и свободное время работников.\\
\bottomrule\end{longtable}\endgroup

\textbf{Дополнительный вклад.} Дополнительный вклад — выполненные выводы и контрольные случаи с независимо обоснованным ответом, датой и результатом выпуска. Уже выполненная эквивалентная проверка используется повторно. Это снимает повторную техническую работу, но само по себе не доказывает нормативное покрытие условия: правовой источник проверяется отдельно. При подтверждении обоих покрытий дополнительная норма не требуется; при доказанно отсутствующем правовом условии необходима его компетентная нормативная фиксация в основном акте или законно предусмотренном инструменте. Общие обязанности и искусственные случаи не доказывают предельную пользу каждого пункта.

\textbf{Соразмерное раскрытие.} Пункт 7 связывает действующую методику и даты на существующем официальном ресурсе; пункт 19 использует существующий канал заявителя для оснований и возражений; пункт 23 хранит протокол внутри и применяет действующую законную отчётность. Отдельный публичный отчёт до выпуска или периодически больше не требуется. При уже обязательной либо законно разрешённой публикации выписка использует внутренний протокол и проверяется на раскрытие. Новая обязанность не обосновывается вымышленным измеренным эффектом публикации.

\textbf{Ресурсы и исполнение.} До внесения 1 декабря разработчики могут согласовать таблицу правовых решений и сопоставить существующую приёмочную документацию с приложенными контролями в рамках действующего задания на проект. До первого юридически значимого применения с 1 января фиксируются действующее правило, требуемые технические/защитные заключения и выполненные проверки затронутых путей. Это последовательность зависимостей, не обещанная длительность и не полномочие отложить внедрение. Дополнительный ручной резерв на работников «Инсон» не возлагается; обычные заявления сохраняют законный процесс. Тесты выполняются в разрешённой среде до внедрения, без дополнительного посещения или документа от семьи.

Для действительно дополнительной задачи разработчик определяет существующий договор либо подразделение, результат, дополнительные часы, законные средства и зависимость завершения. Операционные сведения о сроках указывают период выгрузки, каналы, поступившие и завершённые дела, возраст ожидающих, время обработки и задержки источников; накопительный счётчик сайта этих величин не даёт. Дополнительная обработка каждого заявления вводится только с операционным обоснованием и необходимым испытанием нагрузки, а не утверждением свободного штата. Скорость небольшой авторской программы не выдаётся за пропускную способность YAMIH. Постановление № 687 даёт условные правила стоимости, ПП-388 — бюджетную основу, не доказанное финансирование проекта. Предложение поэтому не утверждает нулевой дополнительный расход и не делает нефинансированный общереспубликанский резерв обязательным обещанием. Точная стоимость интеграции и общая способность Агентства соблюдать сроки не измерены; запрашивается узкая задача использования и выявления пробелов внутри порученного проекта, с отдельной оценкой реального расширения до принятия обязательств.

\section*{Вычислительное задание}
Это рабочее техническое обоснование требований. Для данного проекта коэффициенты и границы официальной многомерной модели 2027 года и данные семей ещё не получены. Отдельно выполнен аудит открытых правил постановления № 35. Общие математические обозначения не являются официальными критериями или названиями категорий.

Утверждённое правило в общем виде \(C=g(F(x),q,v)\): \(x\) — входы, \(q\) — качество и исключения, \(v\) — версия. \(F\) не обязана быть линейной. Действующие основания отказа и назначения помощи учитываются отдельно. Только для аддитивного правила:
\[
S=\sum_jw_jz_j,\qquad z_j\in[a_j,b_j],
\]
\[
L=\sum_j\min(w_ja_j,w_jb_j),\qquad U=\sum_j\max(w_ja_j,w_jb_j).
\]
Если входы независимо принимают любые рациональные значения своих интервалов, граница точна. Для дискретных или зависимых входов это консервативный внешний интервал, способный включать недостижимые категории. Границы требуют методического или информационного основания. При нескольких категориях сохраняется неопределённость; расчёт не даёт права подставить неизвестное значение или отказать.

Тестовая программа использует \(C(S)=|\{t:t\leq S\}|\): равенство переходит к более высокому индексу. Это лишь тестовое соглашение; официальные равенство, округление, исключения и направление отдельно утверждаются.

Повторение: \(R=\text{число расхождений категорий}/\text{число сравнимых случаев}\). Для чувствительности \(F_s=\text{сменившие категорию}/\text{случаи с двумя определёнными категориями}\). Указываются сценарий, числитель, знаменатель, исключённые и неясные случаи. При нулевом знаменателе показатель не определён.

Электронное вычислительное дополнение к материалу 3 содержит аддитивную демонстрацию, точную рациональную арифметику, замкнутые интервалы, равенство на границах и сравнение с контрольной категорией. Отдельно представлены аудит опубликованных правил и независимая арифметическая проверка. Выводы, необходимые для оценки предлагаемых норм, воспроизведены в этой записке; доступ к репозиторию не требуется.

\section*{Выполненный аудит и конкретные предложения}
\textbf{Дата аудита:} 30 сентября 2026 года. Точные входы, контрольные расчёты и результаты исполнения предоставляются в электронном вычислительном дополнении к материалу 3. Основные выводы воспроизведены в следующей таблице. Это аудит опубликованных правил и сконструированных случаев, а не наблюдение реальных семей или проверка официальной модели 2027 года.

\begingroup\fontsize{11bp}{13.2bp}\selectfont
\begin{longtable}{p{.27\linewidth}p{.29\linewidth}p{.35\linewidth}}
\toprule
Проверка & Расчёт & Значение и ограничение \\
\midrule\endhead
Официальный пример, пункт 30 приложения 1 № 35 & $6\,480\,000/5/3=432\,000$; различие ноль & Для сентября 2025 года сохранён май–июль; общий календарный сдвиг уточняется. \\
Пункты 32, 33(a), 39(b); равенство $P$ & При $P=715\,000$ оба условия $D>P$ и $D<P$ ложны & Новый трудоспособный заявитель без иных отказов, прежнего статуса и специальной рекомендации. Операционные решения не изучены; PF-258 также строго ниже $P$. \\
Земля, пункт 27; $(U,B)=(100,90)$ м² & $100-90-\min(20,200)=-10$ м² & $0\leq B\leq U$ выполнено; наличие такого кадастра или исключающего ограничения не установлено. \\
Доход того участка & $-10/100\cdot0{,}035\cdot1\,360\,000\cdot1=-4760$ сумов & Сентябрьское $M$ официально; $k=1$ только допущение. Нулевой предел предлагается, не объявляется действующим. \\
Коэффициенты пола, пункт 23 & $0{,}4M=544\,000$/месяц; на человека при четырёх членах $136\,000$ & Различие вменённого вклада, не измерение заработка, дискриминации или ошибок помощи; поздний $M$ назад не применяется. \\
Расширенная земельная область & $U=100,1000,1001,2000$, все целые $B=0,\ldots,U$: 4105, отрицательных 620 & Сконструированные количества, не доля семей. При $U\leq1000$: $B>0{,}8U$; при $U>1000$: $B>U-200$. \\
Тяжёлая ситуация & Пункты 34–35: $P<D\leq2P$ & Сохраняются рекомендация, квота, обстоятельства, оценка и другие требования; общий интервал через $D$ или $0{,}75D$ не воспроизводит весь путь. \\
Налог и козы, пункты 24, 28 & $\max(2{,}5T/3,B_0)$, переход $T=6B_0/5$; 14 коз: 0, 15: $34\,000$/месяц & Периоды налога и базы согласуются; коэффициент $0{,}07$ и освобождение первой головы. Другие животные и округление не полностью реализованы. \\
\bottomrule\end{longtable}\endgroup

\textbf{1. Согласование границ.} Пункты 5 и 12 предлагаемого приложения включают равенство в обязательные тесты. Окончательное правовое решение согласуется с вышестоящим актом; техническая замена строгого сравнения не предлагается.

Если равенство относится к основной категории, в пунктах 3(a), 3(b) PF-258 и 39(a), 39(b) приложения 1 № 35 условие «доход ниже минимальных потребительских расходов» согласуется с «доход не превышает минимальные потребительские расходы». Примечание (a) приложения 1 PF-258 меняется согласованно.

Для пунктов 3(v) PF-258 и 39(v) № 35 предлагается разделить два доходных пути. Вариант критерия для обсуждения:
\begin{quote}
«К категории семьи на границе бедности относятся:

семьи, выведенные из категорий семьи на государственном обеспечении или бедной семьи, чей месячный средний доход на человека без уменьшения по тяжёлой ситуации выше минимальных потребительских расходов, но не более полутора их размеров;

семьи, включаемые по специальному порядку тяжёлой жизненной ситуации на основании рекомендации “махаллинской семёрки”. На этом пути проверяется непревышение полутора размеров минимальных потребительских расходов после уменьшения дохода на 25 процентов; равенство уменьшенного дохода минимальным расходам либо меньший доход сами по себе не являются основанием отказа. Сохраняются законодательные условия обстоятельств, первичной оценки, рекомендации, предельного количества и повторной оценки».
\end{quote}
Второй путь согласуется с примечанием (b) приложения 1 PF-258 и пунктами 34–35 № 35. Обычный и уменьшенный доход не приравниваются. Вариант требует оценки последствий для права и категории, изменения президентского акта и иных норм; одного постановления недостаточно. Если выбрана иная политика равенства, её категории, отказы и исключения также задаются открыто и согласованно.

Контрпример после рецензии, не входивший в первоначальный список: сконструированная семья из трёх человек с одиноким родителем, законной рекомендацией, квотой и без других отказов; доход за три месяца 8\,580\,000. Обычный $D=2\,860\,000/3>715\,000$; уменьшенный $A=(2\,860\,000/3)\cdot3/4=715\,000$. Верхний предел пункта 35(b) соблюдён. Общая нижняя граница строго выше $P$ противоречила бы этому случаю и не применяется. Двухпутевой вариант его сохраняет. Это проверка условий, не наблюдаемая семья.

\textbf{2. Отрицательная площадь.} Для пункта 27 приложения 1 № 35 предлагается:
\begin{quote}
«$TEM=\max\{0;UM-QM-\min(0{,}2UM;200\text{ м}^2)\}$. Min выбирает наименьшее, max — наибольшее. TEM неотрицательна; единицы UM и QM совпадают. Если данных нет или $QM>UM$, сведения уточняются в компетентном источнике; неизвестное или противоречивое не заменяется нулём. Уточнение само по себе не вводит новое основание отказа».
\end{quote}
Это предложение изменения, не доказательство отсутствия либо наличия такого ограничения в программе. Налоговый коэффициент и другие условия сохраняются. Дата, повторение и влияние на прежние решения определяются окончательным актом.

\textbf{3. Период и версия.} Пункты 4, 6 и 8 фиксируют месячную или периодную природу, календарные месяцы и действующий тогда параметр. Пока период сравнения пункта 26 с $2M$ не выяснен, полная миграционная формула не реализована. Курс по пункту 19 на дату получения YAMIH требует сохранить дату и значение для повторения.

Два условных прочтения не объявлены алгоритмом Агентства. При месячном $q$ переход $2M-1$ к $2M$ для перечисленных стран меняет $3M$ на $2M$, снижение 1\,360\,000 сумов. Для трёхмесячного $q$ меняет $3M$ на $2M/3$, снижение $9\,520\,000/3$. Для прочих стран снижения 0 и $5\,440\,000/3$. Это не сумма помощи. Разработчик определяет и обосновывает период и переключение; программист не добавляет минимальный доход самостоятельно.

Нулевой предел земли может повысить отрицательный нормативный вклад. При сентябрьском $M$ и тестовом $k=1$ максимальное месячное повышение в заявленной области 95\,200; для человека оно делится на состав семьи. Автоматическое увеличение помощи не утверждается. Категории, выплаты и прежние решения требуют законного отдельного повторения. Результаты не обосновывают отмену коэффициентов, сокращение помощи или определённую долю сменивших категорию семей: нужны официальный модельный текст, законные данные и предварительный эмпирический план.

\section*{Сравнение действующих и предлагаемых решений}
Ниже кратко изложен предмет действующей нормы, не её полная дословная копия. Предложения согласуются с окончательным уровнем и редакцией акта.

\textbf{1. Категория и равенство.} Источники: пункт 3 и примечания приложения 1 PF-258, пункты 32–35 и 39 приложения 1 № 35. Предлагается приведённый двухпутевой вариант либо иной согласованный выбор разработчика. Основание: несовпадение признания и категории при $P$ с сохранением специального пути. Вариант не является обязательным условием принятия основного приложения о проверке.

\textbf{2. Земля.} Источник: пункт 27 приложения 1 № 35. Предлагаются нулевой предел и проверяемый источник площадей одного кадастрового объекта. Основание: отрицательная площадь в объявленной алгебраической области. До внедрения оцениваются повышения дохода и прежние выплаты; геометрически обоснованный иной вариант также формулируется воспроизводимо.

\textbf{3. Основания расчёта.} Источник: SMS пункта 37 и соответствующие решения приложения 1 № 35. Дополнительная норма: «Вместе с решением через кабинет предоставляются версия, основные значения и периоды, причины вменения или исключения, правовые и фактические основания и путь возражения. Не использующий кабинет получает сведения через “Инсон”. SMS их не заменяет». Основание: предлагаемые пункты 18–20, статья 24 закона о персональных данных и соответствующие процедуры. Защищённые сведения третьих лиц исключаются.

\textbf{4. Возражение и совет.} Пункты 26 и 31 приложения 4 № 35 ограничивают совет документами «махаллинской семёрки»; пункт 33 устанавливает пятнадцать дней. Новая общая компетенция по автоматике не предполагается. Для форм YAMIH и «Инсон» предлагается: «Возражение против решения только по автоматизированной обработке направляется собственнику или оператору, с письменным результатом в десять дней. Жалоба на махаллинский документ направляется совету по его предмету. Для разных вопросов сохраняются первоначальная дата и маршрутизация». Основание: статья 24 ЗРУ-547. Исправление или дополнение источника по запросу субъекта имеет отдельный трёхдневный срок статьи 11, а выявленные неточные данные исправляются немедленно; этот срок не заменяется десятидневным сроком возражения против автоматизированного решения. Если разработчик расширяет совет, отдельно согласуются предмет, требование доказательств, решение и исполнение; это предложение не утверждает совершённого расширения.

\textbf{5. Сбой.} Пункт 48 приложения 1 № 35 разрешает временно менять алгоритм. Предлагается дополнение: «Записываются основание, область, начало, условие восстановления и расчёты временной версии. После восстановления случаи определяются и проверяются по применимым законным правилам». Основание: не превращать временное полномочие в скрытую постоянную замену критерия; пункты 11, 23 и 26 нового приложения. Полномочие не расширяется. В первый рабочий день публикуются состояние восстановления, доступные результаты, неразрешённые случаи и план; окончательные результаты — при завершении без продления законных индивидуальных сроков.

\textbf{6. Обмен и выплаты.} Согласуются информационные приложения № 35 и назначение и выплаты по приложению 3. В поля передачи добавляются период, время, качество и версия; новая техническая неясность не становится основанием автоматической остановки законной помощи. Исправление дохода, новый критерий политики и взыскание различаются. Конкретные пункты выплат разработчик включает в официальную таблицу после сопоставления выбранной модели 2027 года; охват всех программ этой таблицей сейчас не заявлен.

\textbf{7. Действующая ответственность и дополнительная проверочная процедура.} Пункт 46 приложения 1 к постановлению № 35 уже возлагает на ответственных работников государственного учреждения «Центр информационных технологий» при Агентстве обеспечение правильной работы Реестра и YAMIH, совершенствование модуля и обмен данными; руководители соответствующих министерств и ведомств отвечают за своевременные и достоверные ответы. Пункт 49 возлагает контроль на Генеральную прокуратуру, службу внутреннего аудита Агентства и Социальную инспекцию. Предложение не создаёт новую инспекцию и не передаёт их законные полномочия другому органу. Оно определяет проверяемые материалы для этих обязанностей: независимо воспроизводимые контрольные расчёты, приёмку каждой версии, сохранение периодов и исключений, число неразрешённых случаев, основания решений и допустимые к публикации сводные результаты. Эти два пункта сами не устанавливают такой подробный протокол. Это сопоставление с указанными нормами, не доказательство отсутствия эквивалентных обязанностей или средств во всём законодательстве или действующей системе. Разработчик должен определить существующие эквиваленты, сохранить их полномочия и объединить дублирующиеся нормы. Подтверждённый эквивалент позволяет избежать дублирования, но не оставить отдельное требование без ответа.

\subsection*{Основное проверочное приложение: связь с действующими обязанностями}
Сопоставление относится к обязательной части материала 1. Первая графа излагает предмет конкретной действующей нормы; вторая — дополнительные способы её исполнения, предложенные автором. Это аналитическое сопоставление, а не дословная сравнительная редакция всего постановления. Нумерация во второй графе относится к приложению материала 1.
\begingroup\fontsize{11bp}{13.2bp}\selectfont
\begin{longtable}{p{\dimexpr\linewidth/3-2\tabcolsep\relax} p{\dimexpr\linewidth/3-2\tabcolsep\relax} p{\dimexpr\linewidth/3-2\tabcolsep\relax}}
\toprule
Действующая обязанность & Предлагаемое исполнение & Проверяемая необходимость\\
\midrule\endhead
Пункт 46 приложения 1 к постановлению № 35: правильная работа Реестра и YAMIH в соответствии с законодательством и технической документацией. & Пункты 4–7, 12–17, 25–26: связанное описание правил, независимые ожидаемые ответы, проверка компонентов, приёмка конкретной версии. & Установить соответствие утверждённой норме на предъявленном составе случаев; совпадение итогов или двух программ не скрывает ошибку компонента либо общий неверный ответ.\\
\midrule
Пункты 45–46: защита сведений и своевременные достоверные ответы государственных источников. & Пункты 8–10, 21–22, 28: различимые состояния данных, компетентное уточнение, законный допуск, проверка предоставленного кода и ограничение хранения. & Не превращать отсутствие ответа в ноль, не перелагать государственный источник на семью и не передавать защищённые входы внешнему сервису проверки.\\
\midrule
Пункт 37: SMS о результате. Статья 24 закона о персональных данных: объяснение автоматизированного решения и десятидневное рассмотрение возражения. & Пункты 18–20: доступное объяснение расчёта, разделение источника, программы и правила, сохранение первоначального обращения и компетентного маршрута. & Семья может установить причину и оспорить соответствующее основание; исправление источника не подменяется жалобой на махаллинский документ.\\
\midrule
Пункт 48: временное полномочие при нарушении обмена. Пункт 49: контроль в пределах полномочий перечисленных органов. & Пункты 11, 17, 23–24, 27: запись временного режима, область влияния, восстановление, внутренний протокол и законный пересмотр прежних решений. & Временная версия остаётся обнаружимой; техническое исправление не превращается в самостоятельное основание уменьшения или взыскания выплаты.\\
\bottomrule\end{longtable}\endgroup
Изъятие необязательных вариантов равенства и земельной формулы не исключает эти проверочные обязанности. Подтверждённая равнозначная действующая процедура сохраняется; разработчик указывает её норму и объединяет дублирование. Критерии нуждаемости 2027 года этим сопоставлением не утверждаются.

\subsection*{Трёхграфная таблица по выбранным нормам}
Таблица использует структуру трёх граф пункта 85 приложения к ЗРУ-682. Первая графа содержит только относящуюся к изменению фразу или расчётную часть; остальные условия не отменяются. Это таблица выбранных норм авторского предложения, а не окончательная таблица официального разработчика по всем актам и программам помощи.
\begingroup\fontsize{11bp}{13.2bp}\selectfont
\begin{longtable}{p{\dimexpr\linewidth/3-2\tabcolsep\relax} p{\dimexpr\linewidth/3-2\tabcolsep\relax} p{\dimexpr\linewidth/3-2\tabcolsep\relax}}
\toprule
Соответствующая часть действующей редакции & Предлагаемая редакция & Основание изменения\\
\midrule\endhead
PF-258 пункт 3(a), (b), примечание (a) приложения 1 и постановление № 35, приложение 1, пункт 39(a), (b): условие дохода «\textit{\uline{ниже минимальных потребительских расходов}}». & В выбранном включающем равенство варианте заменить только это условие на «\textbf{не превышающих минимальные потребительские расходы}». Совместно изменить все указанные части на уровне Президента и Кабинета Министров; остальные условия сохранить. & При доходе точно \(P\) у нового заявителя в выбранных условиях не совпадают предикаты признания и основной категории. Вариант является выбором политики и требует отдельной оценки затрат и перехода.\\
\midrule
Постановление № 35, приложение 1, пункт 35(b): специальный путь допускает уменьшенный на 25 процентов доход, не превышающий \(1.5P\), при рекомендации, квоте и других условиях. & \textbf{Разделить путь обычного дохода и специальный путь уменьшенного дохода. На специальном пути проверять \(A\leq1.5P\); само по себе \(A\leq P\) не становится основанием отказа. Остальные законные условия сохранить.} & Новый общий нижний предел \(A>P\) противоречит \(A=P\) в сконструированном контрольном случае семьи из трёх человек. Общий интервал не должен сокращать специальный путь.\\
\midrule
Алгебраическая запись соответствующего расчёта площади в пункте 27 приложения 1 № 35: \textit{\uline{\(L=U-B-\min(U/5,200)\)}}. & После подтверждения геометрического смысла разработчиком предлагается \textbf{L = max\{0, U − B − min(U/5, 200)\}}. Площади относятся к одному объекту и периоду; нулевой предел, источник и единица определяются нормативным правилом. & При \(0\leq B\leq U\) возможен отрицательный результат. Нулевой предел может повысить вменённый доход; он не означает автоматического увеличения помощи или взыскания прежних выплат. Операционная частота неизвестна.\\
\midrule
Пункт 37 приложения 1 № 35: заявителю на мобильный номер направляется SMS о включении семьи в реестр или об отказе. & Сохранить SMS. \textbf{Вместе с решением открыть в кабинете версию, основные значения и периоды, правовые и фактические основания и путь возражения; не использующий кабинет получает сведения через «Инсон». Защищённые сведения третьих лиц исключаются.} & Сообщение результата не раскрывает полностью расчёт и порядок возражения. Дополнение согласуется со статьёй 24 закона о персональных данных и применимыми административными процедурами.\\
\bottomrule\end{longtable}\endgroup

\section*{Обсуждение, согласование и основания ссылок}
Авторское предложение не направлено государственным органам, не проходило общественного обсуждения на государственном портале и официального межведомственного согласования. Независимая внутренняя проверка подтверждает качество расчётов в указанном объёме и не является межведомственным согласованием. Доказательств согласия или возражений ведомств нет; ответы на возражения рассматривают возможные вопросы, выделенные автором.

Статья 22 обосновывает участие в подготовке, научные рекомендации и консультации специалистов в основном пункте 4. PF-193 определяет поручение и предмет. PF-258 и постановление № 35 определяют действующие правила категорий, расчёта, обращений и временных полномочий. Эти ссылки не разрешают нижестоящему акту или технической инструкции заменить критерий большей юридической силы. Законы о персональных данных и административных процедурах обосновывают контроль данных, доступ к основаниям решения и законный путь возражения.

\section*{Последовательность исполнения и ответственность}
После компетентного утверждения методики существующие методические, системные и юридические функции согласуют единицы, периоды, границы, исключения, даты и основания. Независимые случаи включаются в применимый технический процесс и выпуск; равноценные тесты указываются по ссылке. Приёмка пунктов 25–26 не заменяет обязательные внешние заключения. Описание связывается через существующий ресурс, протокол хранится внутри по пункту 23. Публикация следует действующей обязанности или разрешению; новый публичный отчёт не условие запуска. Сроки и январское начало сохраняются.

Компетентный собственник или оператор решает вопрос фактических данных, системное подразделение — программный дефект, методическое подразделение с юридической службой готовят правовой анализ. При необходимости официального толкования или изменения вопрос направляется компетентному принявшему акт органу; статья 55 не позволяет толкованием вносить изменения. Индивидуальное решение меняет компетентный субъект. Приёмка не утверждает критерии политики. Поручение и законные сроки не переносятся внутренним планом.

\section*{Исполнение по этапам и доказательства завершения}
\begingroup\fontsize{11bp}{13.2bp}\selectfont
\begin{longtable}{p{0.22\linewidth} p{0.25\linewidth} p{\dimexpr0.53\linewidth-6\tabcolsep\relax}}
\toprule
Задача и пункты & Ответственный & Материал и условие завершения\\
\midrule\endhead
Методика и изменение; 4–7, 12 & Агентство с Министерством; методическое подразделение и юридическая служба & Паспорта правил и матрица изменений. Правовые основания, даты, приоритеты и все применимые пути согласованы; отсутствующее правило не вводится программой.\\
\midrule
Источники и качество; 8–10 & Ответственные работники Центра информационных технологий; руководители предоставляющих сведения ведомств в пределах пункта 46 & Карта источников и журнал качества. Для каждого использованного поля установлен период, источник, состояние и способ уточнения; ноль отличается от неизвестного.\\
\midrule
Контрольная реализация; 12–17 & Системное подразделение и проверяющий контрольные ответы, отличный от ответственного за реализацию & Реестр правил и случаев, точные контрольные ответы и протокол. Состав случаев совпадает с исполненным; все влияющие расхождения устранены и повторно проверены.\\
\midrule
Объяснение и возражение; 18–20 & Уполномоченный работник «Инсон», собственник или оператор данных, орган пересмотра в его компетенции & Образец доступного объяснения и связанная запись направления. Получатель, срок и результат установлены; проверка маршрута не задерживает законный срок.\\
\midrule
Сбой и восстановление; 11, 23 & Уполномоченный и ответственные системы & Существующие записи временной версии и восстановленных расчётов; решения выявляются, различия разрешаются законно; внутренние доказательства обновляются.\\
\midrule
Ресурсы и приёмка; основные 6–7, 25–26 & Агентство с Министерством; руководитель или уполномоченный заместитель & Запись выпуска, внешние заключения, тесты, реальные дополнительные задачи/средства. Мощность новой ручной службы не предполагается.\\
\midrule
Хранение и законное раскрытие; 21–24, 28 & Оператор, защита и действующий надзор & Законный график, внутренний протокол и проверка обязательного/разрешённого раскрытия; нового цикла публикации нет.\\
\bottomrule\end{longtable}\endgroup

До внесения основного проекта 1 декабря 2026 года согласуются законная таблица решений и связь контролей с существующими приёмочными доказательствами. До первого применения затронутого правила с 1 января 2027 года фиксируются компетентное утверждение, применимые технические/защитные заключения и проверки путей. Это последовательность зависимостей, не оценённый кадровый график и не полномочие отложить начало. Предварительные случаи готовятся по проекту; окончательный ответ следует утверждённому правилу. Дополнительная интеграция требует реальных сведений о задаче и средствах до обязательства расходов.

\section*{Предварительное воздействие и финансовое обоснование}
\textbf{A0 — действующий порядок.} Контроли и законные права сохраняются. Их правовое существование установлено; доказательства выпусков Агентства для определения охвата случаев не получены. Это не вывод о неработающем контроле.

\textbf{A1 — техническое исполнение правовой основы.} Выводы и случаи включаются в существующие требования и записи выпуска при отдельно установленном законном правиле. Равноценные проверки и каналы используются повторно; непокрытая задача и ресурсы документируются. A1 — предпочтительный способ технического исполнения, а не альтернатива отсутствующей нормативной основе.

\textbf{A2 — минимальное нормативное содержание порученного акта.} В проект по PF-193 включаются необходимые условия новой методики и не обеспеченные конкретной действующей нормой гарантии. Можно выбрать пять кратких пунктов материала 1, равноценную редакцию либо сохранённую применимую норму; A1 обеспечивает техническое исполнение. Полная форма не неделима. Не предлагаются отдельная ручная служба, новая статистическая отчётность или переход Президента по правам. Сравнение не является официальным отчётом по статье 35 ЗРУ-682 и пункту 2(b) ПП-5025.

Оцениваются только выявленные дополнительные задачи: правовое согласование, недостающие случаи, затронутые изменения интеграции/доступа и фактическое сопровождение. Существующая работа и покрытие договором фиксируются против повторной закупки и двойного учёта. Агентство предоставляет объёмы, часы, операционные доказательства и договорную область; Министерство — законные средства. Интеграция Агентства, избыток персонала, итоговая стоимость и экономия не измерены. Существующая программа и код не доказывают нулевую стоимость.

Эмпирическая нуждаемость требует отдельного исследования на законных данных. Его стоимость и результат можно отделить от проверки соответствия; незавершённое исследование не доказывает правильную адресность.

\section*{Доказательства приёмки и финансирования}
Приложение проверки может приниматься отдельно от необязательных поправок равенства и земельной формулы. Оно требует реестра связей правил и случаев, независимо проверенного правового вывода ожидаемых ответов и компонентного сравнения. Нельзя пропустить применимое исключение и затем удостоверить полноту собственного списка. Отсутствующие или повторные случаи, неразрешённые компонентные различия и нулевые знаменатели раскрываются. Совпадение сумм не скрывает компенсирующие ошибки. Методика, программный выпуск и датированная конфигурация имеют отдельно различимые записи, позволяя восстановить программное исправление при неизменном правиле.

Выполненная работа предоставляет 37 регрессионных тестов и проверяющий фиксированный охват случаев. Она не предоставляет производственную приёмку: таблица 5 научного обоснования для каждой группы правил указывает выполненные проверки и необходимые дополнительные доказательства официального источника, ветвей и округления. Десять основных строк дохода и 4\,105 земельных строк являются построенными контролями. Они не считаются семьями и не используются как прогноз бюджетного эффекта.

Протокол готовности разработчика указывает методическое подразделение, подразделение системы и юридическую службу, принятые материалы, применимые даты и неразрешённые вопросы. Учёт ресурсов фиксирует объёмы, дополнительные денежные затраты, часы, доступную мощность и финансирование интеграции, контроля, хранения, доступа, обучения и пересмотра. Постоянные и переменные расходы не считаются дважды. Итог бюджета либо экономия не выдуманы. Последующее исследование благосостояния можно оценить отдельно; дефицит мощности не продлевает законные сроки заявителя.

Технические источники устройства контроля: \href{https://discovery.ucl.ac.uk/id/eprint/1471263/1/06963470.pdf}{Barr и соавторы, 2015, обзор проблемы ответа теста} и \href{https://nvlpubs.nist.gov/nistpubs/ir/2021/NIST.IR.8397.pdf}{NIST IR 8397, разделы 2.6–2.8}. Это инженерные источники, не дополнительные правовые обязанности Узбекистана.

\section*{Ожидаемые результаты, последствия и проверяемые показатели}
Ожидаемый непосредственный результат — наличие у каждой применённой версии описанного правила, проверенного контрольного комплекта и сохранённого основания индивидуального решения. Исполнение оценивается по наличию паспортов всех применимых критериев, совпадению заданного и выполненного состава тестов, числу неразрешённых расхождений компонентов, доступности объяснения и соблюдению законных сроков исправления и возражения. Для долей публикуются числитель, знаменатель и исключения. Эти показатели проверяют выполнение порядка; они не являются показателями точности определения нуждаемости. Проект не устанавливает произвольный процент допустимых ошибок, устойчивости или сокращения жалоб.

Для семьи ожидаются доступ к использованным значениям и правовому основанию, возможность исправления сведений у компетентного источника и отсутствие нового отказа только из-за технического статуса. Для органов появляются обязанности вести связанный контрольный комплект, готовить объяснение и отчёт, обеспечивать необходимую мощность. Устранение ошибки расчёта может изменить законный результат в любую сторону; изменение выплаты, пересмотр и взыскание требуют собственных правовых оснований и компетентного индивидуального решения. Действующие самостоятельные основания отказа, исключения из Реестра и прекращения выплаты сохраняются. По синтетическим случаям число таких изменений и их бюджетная стоимость не определяются.

Практические риски — неполный источник, ошибочный контрольный ответ, неподготовленный маршрут обращения, нехватка работников и раскрытие через малые группы. Меры проекта соответствуют рискам: раздельный статус и уточнение (8–10), независимый правовой вывод и проверка компонентов (12–17), объяснение и маршрутизация (18–20), защищённый допуск и публикация (21–23), ресурсный расчёт и письменная приёмка (основные 6–7, 25–26). Наличие контрольной программы само по себе не устраняет эти риски; условия закрытия указаны в таблице исполнения.

\section*{Состав финансово-экономического расчёта}
К основному проекту предлагается подготовить расчёт по следующим работам: описание и юридическое согласование правил; интеграция источников и полей качества; независимые контрольные ответы и выполнение тестов; журнал версий и защищённое хранение; кабинет и обслуживание обращений через «Инсон»; обучение работников; отчёт и контроль раскрытия; сопровождение и повторные проверки изменений. По каждой работе отдельно записываются разовый объём, годовой объём, единица, дополнительные денежные затраты на единицу, часы работников, действующий ресурс, требуемое увеличение и документ об источнике финансирования. Наличие работника в штате не означает свободной мощности; перераспределённый труд показывается отдельно от новых денежных затрат.

Используется тот же учёт, что в материале 3:
\[K_{\rm first}=K_{\rm setup}+K_{\rm annual},\qquad K_{\rm annual}=K_{\rm fixed}+\sum_r Q_r k_r.\]
Здесь \(Q_r\) — годовой объём работы в указанной единице, \(k_r\) — дополнительные денежные затраты на единицу; фиксированная и переменная части не дублируются. Агентство определяет объём и мощность по собственным данным, Министерство — законный источник и бюджетное оформление. Средства на независимое изучение фактической нуждаемости и на необязательные изменения границ рассчитываются отдельно. Итоговая стоимость проекта не определена: утверждённый состав работ, число модулей, длительность и выделенное финансирование не получены; опубликованный порядок оценки стоимости указан ниже. Поэтому не заявляется ни нулевая стоимость, ни установленная сумма дополнительного финансирования; приведён состав необходимых подтверждений для основного пункта 6.

\subsection*{Проверенные ведомственные положения об оценке стоимости}
В действующей редакции \href{https://lex.uz/docs/-6720230}{постановления Кабинета Министров от 28 декабря 2023 года № 687}, основной пункт 3 и приложение 3, установлен порядок определения стоимости проектов Центра информационных технологий Агентства за счёт Государственного фонда социальной защиты либо внебюджетных средств Агентства. Применение зависит от исполнителя и источника средств; это не общий тариф на любую проверочную работу.

Для технического сопровождения и совершенствования информационных систем пункт 4 приложения 3 определяет фонд оплаты труда:
\[\mathrm{MHTF}=1{,}85\times\mathrm{MS}\times16\times\mathrm{MHTEKM}\times\mathrm{LAOOS}.\]
MS — число модулей, LAOOS — число месяцев проекта, MHTEKM — установленный минимальный размер оплаты труда с увеличением на социальный налог, если Центр является его плательщиком. Иные связанные с проектом расходы определены в размере 5 процентов фонда оплаты труда. Пункты 3 и 7–9 отдельно предусматривают налоги и обязательные платежи, содержание Центра и управленческого/вспомогательного персонала и прибыль; формула фонда не является полной ценой договора. Для проектов пунктов 4–6 пункт 8 ограничивает содержание Центра и управленческого/вспомогательного персонала суммой до 40 процентов стоимости проекта, определённой по этим пунктам; число работников и условия оплаты утверждаются приказом директора Агентства. По пункту 9 прибыль Центра составляет 2 процента расходов по проекту. Проценты применяют к указанной базе и виду проекта, не без разбора ко всем задачам. Для новых проектов пункт 10 предусматривает определение стоимости по техническим требованиям, сроку, категориям и числу специалистов на основании приказа директора Агентства. По пункту 12 условные штатные и зарплатные параметры используются для оценки стоимости, а не доказывают фактический штат либо свободную мощность. Перед расчётом определяется вид работы; число модулей и длительность здесь не предполагаются.

\href{https://lex.uz/docs/-7959569}{ПП-388 от 26 декабря 2025 года}, пункты 19–21 и приложение 6, предусматривает ведомственное программное бюджетирование. Это бюджетный контекст, а не выделение средств на данный контрольный модуль. Агентство и Министерство связывают итоговый расчёт с применимой утверждённой программой, сметой и решением о финансировании. Общие расходы на социальную защиту не подменяют цену предлагаемой проверки расчётов.

\section*{Международные материалы и языки комплекта}
Иностранная норма не заимствуется и новый международный договор не исполняется этим проектом. OECD/JRC, обзор контрольных ответов Barr и рекомендации NIST используются как инженерные источники проверки неопределённости и реализации; их применение к нормативному расчёту является обоснованной в материале 3 адаптацией. Они не устанавливают обязательные для Узбекистана коэффициенты или пороги. Для данного подзаконного предложения сведения о зарубежных материалах раскрываются в объёме, необходимом по пункту 84 Единой методики; иностранная законодательная модель не объявляется доказанно оптимальной.

Основной язык авторского обращения и материалов 1–3 — русский. Статья 6 \href{https://lex.uz/docs/-3336169}{закона об обращениях} допускает подачу на государственном и других языках. Узбекская редакция предоставлена для подготовки официального проекта на государственном языке по статье 21 \href{https://lex.uz/docs/-5378966}{ЗРУ-682}; английская является согласованной сопроводительной версией. При согласовании все три текста сверяются по формулам, операторам неравенства, срокам и полномочиям. Содержательное расхождение переводов устраняется до подачи разработчиком окончательной редакции, а не разрешается произвольным выбором удобного текста.

\section*{Ответы на возражения}
Для данного проекта официальная спецификация баллов 2027 года не получена, поэтому точность и бюджетные последствия новой модели не рассчитаны. Это не препятствует аудиту открытых формул и обязанности будущей проверки. Если отрицательный пример не встречается в кадастре, паспорт указывает доказанную область; реальная частота не заявлена. Если специальный приоритет уже разрешает тяжёлый путь, фиксируются основание и контрольный случай. Конфиденциальность обеспечивается внутренним законным расчётом и контролем выпуска, не передачей автору записей. Ресурсы рассматриваются настоящим расчётом. Уже равноценную действующую норму разработчик указывает и объединяет с дублирующим предложением.

\section*{Процедурные гарантии и пределы подтверждения}
\textbf{Имеющиеся полномочия и необходимость.} Проверяются отдельно правовое покрытие и выполненный контроль. Статья 20 ЗРУ-682 ограничивает инициирование указанных в ней проектов, если вопрос решается существующими законными механизмами и административными полномочиями; PF-193 прямо поручает нормативную разработку новой оценки. Неопределённое правовое условие нельзя закрыть протоколом теста. При равноценном действующем правиле оно сохраняется, а техническая интеграция использует существующие полномочия. Краткие пять пунктов материала 1 либо эквивалентная редакция позволяют нормативный результат без полного приложения. Конкретные авторские пункты не объявляются автоматически обязательными.

\textbf{Приём обращения.} Статья 40 \href{https://lex.uz/docs/-6445145}{Конституции} защищает предложения компетентным органам и требует их рассмотрения в законном порядке и сроки. Статьи 6, 18 и 23 \href{https://lex.uz/docs/-3336169}{закона об обращениях} разграничивают реквизиты, приём и регистрацию. Регистрация производится в день поступления, при поступлении после рабочего времени — на следующий рабочий день; отказ в регистрации запрещён. Обращение указывает автора, место жительства, получателя и вопросы. Статьи 29–30 сохраняют определённые основания оставления без рассмотрения или прекращения рассмотрения, включая анонимность или иное несоблюдение закона, отзыв и установленные условия повторного обращения. Неизвестные обстоятельства будущей подачи или прежней внешней переписки данным обзором не подтверждаются.

\textbf{Рассмотрение и ответ.} Статья 23 ЗРУ-682 требует изучать относящиеся к предмету исследования и обращения, обобщать и использовать предложения и научные рекомендации. Статья 22 допускает участие, но не обязывает назначать. По статье 27 закона об обращениях рассматриваются все вопросы и письменный или электронный ответ содержит конкретные основания. Запрошенный ответ по каждому пункту делает эту обязанность проверяемой; закон не предписывает обязательную табличную форму. Статья 28 устанавливает для предложения срок до месяца, кроме дополнительного изучения с письменным уведомлением в десятидневный срок. Личная просьба об участии или доступе может требовать отдельной классификации и соответствующего срока. Объединённое письмо не устанавливает обязательного судебного толкования этой классификации; исключению для дополнительного изучения не приписан выдуманный конечный срок.

\textbf{Доказательства исполнения.} Сохраняются поданная редакция, приложения, квитанция либо электронное подтверждение, номер регистрации и дата поступления. Сохраняются уведомления о направлении и дополнительном изучении, запросы информации и полный ответ. Ответ сопоставляется с пятью нумерованными просьбами и каждым содержательным модулем: включение, изменённое решение, точная действующая норма либо причина непринятия. Хеш исходного файла показывает, какой файл проверен; он не доказывает получение органом. Законно принятый акт, официальное решение о назначении доказывают соответствующие результаты. Компиляция проектов и успешные тесты программы сами по себе этих результатов не создают.

\textbf{Защита от незаконного отказа.} Статьи 32–33 закона об обращениях предусматривают обжалование незаконного отказа в приёме или рассмотрении и разъяснение порядка обжалования; \href{https://lex.uz/docs/-6445145}{статья 55 Конституции} защищает судебное обжалование незаконных решений, действий и бездействия. Статьи 185–189 \href{https://lex.uz/docs/-3527353}{Кодекса об административном судопроизводстве} задают применимую основу. Статья 189 требует одновременно незаконности и нарушения защищённых прав или интересов заявителя для указанного средства защиты; суд может обязать принять законное решение, совершить действие или иначе устранить нарушение. Несогласие с обоснованным выбором политики само по себе этого нарушения не доказывает. Статья 186 устанавливает общий шестимесячный срок обращения в суд со дня, когда стало известно о нарушении, если законом не предусмотрен иной срок; восстановление требует уважительной причины. Это отличается от годичного срока жалобы вышестоящему органу по статье 22 закона об обращениях. При реальном обжаловании проверяются факты, классификация, компетентный суд и действующие правила. Решение суда, обязывающее принять это предложение, не получено и не выявлено.

\textbf{Пределы правового вывода.} Статья 23 ЗРУ-682 делает предложения обсуждения рекомендательными с обязательным рассмотрением; статья 24 прямо разрешает отклонение с обоснованием причин. Здесь доказательствами поддерживается наличие указанных процедурных обязанностей с законными условиями и средствами защиты. Это не гарантия назначения, доступа к ограниченным сведениям, включения или принятия. Статья 3 \href{https://lex.uz/docs/-3492199}{ЗРУ-457} исключает подготовку и принятие нормативных актов из своего действия. Поэтому требования к индивидуальному административному решению, приведённые в иных частях пакета, не заявляются как процедурная основа принуждения к изданию акта; предложенный порядок обращений опирается на обращение гражданина и отдельные конституционные и законные обязанности, описанные выше.

\textbf{Применимая процедура и судебное доказывание.} Статья 1 закона об обращениях исключает обращения, порядок рассмотрения которых установлен иным законодательством. Данное письмо подготовлено как инициативное предложение гражданина и личная просьба, а не состоявшееся обращение в рамках формального обсуждения. Если орган применяет специальный порядок, в запрошенном ответе следует указать его и правовое основание; обычный срок предложения не заявляется для последующего портального комментария по статье 24 ЗРУ-682. Обязательное решение, устанавливающее классификацию этого смешанного письма, не выявлено. Это предел толкования, не доказательство отсутствия конституционного права автора направить предложение.

При оспаривании в применимом порядке административного судопроизводства \href{https://lex.uz/docs/-3527353}{статья 67 Кодекса} возлагает бремя доказывания на орган или должностное лицо, чьи решение или действие оспариваются, требует подтверждения фактов возражений и ограничивает их основаниями, на которые они опирались при решении или действии. Заявитель участвует в сборе доказательств в пределах своих возможностей и доказывает размер заявленного ущерба. \href{https://lex.uz/docs/-4711311}{Пункты 17 и 20 постановления Пленума Верховного суда № 24} разъясняют бремя доказывания законности и устанавливают, что суд не оценивает целесообразность решения или действия, законно оставленного административному усмотрению. Пункт 21 конкретно разъясняет защиту при нерассмотрении письменного обращения в установленный срок: суд признаёт бездействие незаконным и обязывает рассмотреть обращение в установленный судом срок, не предрешая, какое решение должно быть принято. Эти правила поддерживают проверку действительных оснований отказа и исполнения применимых обязанностей; они не обеспечивают автоматической победы или обязанности суда выбрать авторскую политику.

\section*{Источники и условия завершения}
Поручение: \href{https://lex.uz/docs/-8472526}{PF-193, пункты 13–14}. Формулы и процедуры: \href{https://lex.uz/docs/-8027513}{постановление № 35 с приложениями}. Категории: \href{https://lex.uz/docs/-7952180}{PF-258, пункт 3}. $M$: \href{https://lex.uz/docs/-8283656}{PF-115}; $P$: \href{https://stat.uz/uz/matbuot-markazi/qo-mita-yangiliklar/67179-minimal-istemol-kharazhatlari-ijmati-t-risida-2}{сообщение статистики 4 марта 2026 года}. Участие: \href{https://lex.uz/docs/-5378966}{ЗРУ-682, статья 22}. Данные и автоматические решения: \href{https://lex.uz/docs/-4396419}{ЗРУ-547, статьи 11, 16–18, 24 и 28; защита специальных данных}. Основания индивидуального акта: \href{https://lex.uz/docs/-3492199}{ЗРУ-457, статья 54}. Чувствительность: \href{https://www.oecd.org/content/dam/oecd/en/publications/reports/2008/08/handbook-on-constructing-composite-indicators-methodology-and-user-guide_g1gh9301/9789264043466-en.pdf}{руководство OECD/JRC, 2008, шаг 7}.

Дополнительно: \href{https://lex.uz/docs/5331933}{ПП-5025, пункт 2(b)}; \href{https://lex.uz/docs/-5378966}{методика ЗРУ-682, пункты 11, 43–51, 83–85}. Аналитические тексты — 11–14 пунктов, нормативное тело — 14; A4, поля 3 см слева и 2 см остальные. Liberation Serif — объявленная замена Times New Roman; электронный формат подтверждает разработчик. Официальный класс LaTeX не заявляется.

Это авторское предложение; новые обязанности не объявлены действующим правом. Нужны официальная методика, единая таблица всех выбранных изменений, стоимость, реквизиты применения и человеческая проверка узбекского юридического языка. Модуль не исполняет весь PF-193, особенно все правила направления к услугам. Принятие и отсутствие любых правовых возражений не гарантируются.
\end{document}
